%==========================================================================
%  TRES0312 Fysikk — Foiler-mal
%
%  MØrk/lys modus: kommenter ut én av de to linjene nedenfor.
%==========================================================================
\newif\ifdarkmode
\darkmodetrue      % ← mørk modus
%\darkmodefalse   % ← lys modus  (fjern % foran og legg % foran linjen over)

\documentclass[aspectratio=169, 12pt]{beamer}

\usepackage{amd-beamer}
\usetikzlibrary{overlay-beamer-styles}
\usetikzlibrary{calc}

%--------------------------------------------------------------------------
%  DOKUMENTINFO
%--------------------------------------------------------------------------
\title{Bevaring av mekanisk energi}
\subtitle{TRES0312 --- kapittel 5.4}
\author{Ben David Normann}
\date{\today}

%--------------------------------------------------------------------------
%  SEKSJONSKONTEKST
%--------------------------------------------------------------------------
\setcounter{section}{5}   % Kapittel 5: "Arbeid og energi"

\begin{document}

%--------------------------------------------------------------------------
%  SLIDE 1 — TITTELSIDE
%--------------------------------------------------------------------------
\begin{frame}[plain]
  \titlepage
  \dekkerdelkap{kapittel 5.4}
  \vspace{0.3cm}
  \begin{center}
  {\bfseries\color{repcol} Første time: eksempler på bevaring av mekanisk energi. Etter pausen: loop-the-loop.}
  \end{center}
\end{frame}

%--------------------------------------------------------------------------
%  SLIDE 2 — Tre kjappe fra forrige gang
%--------------------------------------------------------------------------
\begin{frame}{}

  \repetisjon{Tre kjappe fra forrige gang}{%
    \begin{enumerate}[1.]
      \item Hva er den mekaniske energien til et legeme, $E_{\text{mekanisk}} = \ldots$?
      \item Under hvilken betingelse er den mekaniske energien bevart?
      \item Hva sier bevaringsloven, $E_{\text{før}} = \ldots$?
    \end{enumerate}
  }

\end{frame}

%--------------------------------------------------------------------------
%  SLIDE 3 — Repetisjon av definisjonen
%--------------------------------------------------------------------------
\begin{frame}{Bevaring av mekanisk energi}

  \defnat{5.6}{Bevaring av mekanisk energi}{
  Når et legeme kun påvirkes av tyngdekraft og/eller fjærkraft:
  \[
    E_{\text{før}} = E_{\text{etter}}, \qquad E = m\mathrm{g}h + \frac{1}{2}mv^2 + \frac{1}{2}kx^2
  \]
  }

  \pause
  \oppsum{Algoritme for løsning}{%
  \begin{enumerate}[1.]
    \item[] Velg nullnivå for høyden $h$.
    \item[] Skriv opp $E_{\text{før}} = E_{\text{etter}}$ med alle energiledd på hver side.
    \item[] Stryk fellesledd (ofte $m$), og løs for den ukjente.
  \end{enumerate}
  }

\end{frame}

%--------------------------------------------------------------------------
%  SLIDE 4 — Eksempel: Fallende stein
%--------------------------------------------------------------------------
\begin{frame}{Eksempel: fallende stein}

  \exmat{5.4}{Fallende stein}{%
  En person holder en stein $1.60\,\mathrm{m}$ over bakken og slipper den. Luftmotstand neglisjeres.

  Finn steinens fart rett før den treffer bakken.
  }

  \pause
  Med $h=0$ ved bakken, $v_1=0$ ved start:
  \[
    m\mathrm{g}h_1 = \frac{1}{2}mv_2^2 \quad\Longrightarrow\quad v_2=\sqrt{2\mathrm{g}h_1}
  \]
  \[
    v_2 = \sqrt{2\cdot 9.8\cdot 1.60} = \doubleunderline{5.6\,\mathrm{m/s}}
  \]

\end{frame}

%--------------------------------------------------------------------------
%  SLIDE 5 — Eksempel: Fjærkanon
%--------------------------------------------------------------------------
\begin{frame}{Eksempel: fjærkanon}

  \exmat{5.5}{Fjærkanon i tyngdefeltet}{%
  En fjærkanon skyter en kanonkule ($m=0.2\,\mathrm{kg}$) rett opp etter at fjæren er komprimert $x=0.15\,\mathrm{m}$. Kula når $h=5.0\,\mathrm{m}$.

  Finn fjærkonstanten $k$.
  }

  \pause
  All elastisk potensiell energi blir gravitasjonspotensiell energi på toppen:
  \[
    \frac{1}{2}kx^2 = m\mathrm{g}h \quad\Longrightarrow\quad k = \frac{2m\mathrm{g}h}{x^2} = \doubleunderline{871\,\mathrm{N/m}}
  \]

\end{frame}

%--------------------------------------------------------------------------
%  SLIDE 6 — Eksempel: kollisjon mellom legeme og fjær
%--------------------------------------------------------------------------
\begin{frame}{Eksempel: legeme kolliderer med fjær}

  \exmat{5.6}{Kollisjon mellom legeme og fjær}{%
  Et legeme ($m=100\,\mathrm{g}$) glir friksjonsfritt fra ro, $60\,\mathrm{cm}$ over en fjær med $k=300\,\mathrm{N/m}$.
  \begin{enumerate}[a)]
    \item Finn farten idet legemet treffer fjæren.
    \item Finn maksimal kompresjon av fjæren.
  \end{enumerate}
  }

  \pause
  a) $m\mathrm{g}h=\tfrac{1}{2}mv^2 \Rightarrow v=\sqrt{2\mathrm{g}h}=\doubleunderline{3.43\,\mathrm{m/s}}$

  \pause
  b) $\tfrac{1}{2}mv^2=\tfrac{1}{2}kx^2 \Rightarrow x=\sqrt{\dfrac{mv^2}{k}}=\doubleunderline{0.063\,\mathrm{m}}$

\end{frame}

%--------------------------------------------------------------------------
%  SLIDE 7 — Aktivitet
%--------------------------------------------------------------------------
\begin{frame}{Aktivitet}

  \aktivitet{Pendel}{%
  En pendel med lengde $L=1.2\,\mathrm{m}$ slippes fra ro med snora horisontal.

  Bruk bevaring av mekanisk energi til å finne farten til pendellodet når snora er loddrett (i det laveste punktet).
  }

\end{frame}

%--------------------------------------------------------------------------
%  SLIDE 8 — Oppsummering
%--------------------------------------------------------------------------
\begin{frame}{}

  \oppsum{Bevaring av mekanisk energi}{%
  \begin{enumerate}[1.]
    \item[] $E_{\text{før}}=E_{\text{etter}}$, der $E=m\mathrm{g}h+\tfrac{1}{2}mv^2+\tfrac{1}{2}kx^2$.
    \item[] Gjelder kun når friksjon og luftmotstand er neglisjerbare.
    \item[] Velg nullnivå for $h$ med omhu --- det forenkler ofte regningen.
  \end{enumerate}
  }

\end{frame}

%--------------------------------------------------------------------------
%  SLIDE 9 — PAUSE
%--------------------------------------------------------------------------
\begin{frame}
\begin{center}
  \Huge{Pause}\\[8pt]
  {\large Etter pausen: loop-the-loop}
\end{center}
\end{frame}

%--------------------------------------------------------------------------
%  SLIDE 10 — Til ettertanke: loop-the-loop
%--------------------------------------------------------------------------
\begin{frame}{}

  \tanke{}{%
  En ball ruller ned en bane og inn i en sirkulær loop.\\[4pt]
  Er normalkraften på ballen like stor i bunnen av loopen som på toppen? Hva avgjør om ballen holder kontakt med skinna hele veien rundt?
  }

\end{frame}

%--------------------------------------------------------------------------
%  SLIDE 11 — Eksempel: Loop-the-loop, oppsett
%--------------------------------------------------------------------------
\begin{frame}{Eksempel: loop-the-loop}

  \exmat{5.7}{Loop-the-loop}{%
  En ball med masse $m=1.0\,\mathrm{kg}$ slippes fra ro fra høyde $H=10.0\,\mathrm{m}$ over bunnen av en friksjonsfri sirkulær loop med radius $R=2.0\,\mathrm{m}$.

  Finn normalkraften i bunnen ($h=0$), på siden ($h=R$) og på toppen ($h=2R$).
  }

\end{frame}

%--------------------------------------------------------------------------
%  SLIDE 12 — Figur: loop-the-loop
%--------------------------------------------------------------------------
\begin{frame}{}

  \begin{figure}[ht!]
  \centering
  \resizebox{0.85\textwidth}{!}{%
  \begin{tikzpicture}[scale=0.90, >={Stealth[length=5pt,width=4pt]}]
    \pgfmathsetmacro{\R}{1.6}
    \draw[gray!50, thick] (-3.6, 0) -- (2.9, 0);
    \draw[thick] (-2.8, 4*\R) -- (0.02, 0.02);
    \draw[thick] (0, \R) circle (\R);
    \filldraw[headCol!50!white] (-2.8, 4*\R) circle (0.13);
    \node[right=2pt, font=\footnotesize] at (-2.65, 4*\R) {start, $v_0 = 0$};
    \draw[<->, thin, gray!70] (-3.4, 0) -- (-3.4, 4*\R)
        node[midway, left, font=\footnotesize] {$H$};
    \draw[dashed, thin, gray!40] (-2.8, 4*\R) -- (-3.4, 4*\R);
    \draw[dashed, thin, gray!40] (0, 0) -- (-3.4, 0);
    \draw[<->, thin, gray!70] (2.4, 0) -- (2.4, 2*\R)
        node[midway, right, font=\footnotesize] {$2R$};
    \draw[dashed, thin, gray!40] (0, 2*\R) -- (2.4, 2*\R);
    \fill[gray!50] (0, \R) circle (0.055);
    \draw[<->, thin, gray!70] (0.06, \R) -- (\R, \R)
        node[midway, above, font=\footnotesize] {$R$};
    \filldraw[headCol] (0, 0) circle (0.13);
    \node[below right=1pt, font=\small, headCol!60!black] at (0.1, -0.08) {\textbf{1}};
    \draw[->, very thick, defscol]
        (0.18, 0) -- (0.18, 1.1) node[right, font=\footnotesize] {$N_1$};
    \draw[->, very thick, red!70!black]
        (-0.18, 0) -- (-0.18, -0.85) node[left, font=\footnotesize] {$m\mathrm{g}$};
    \filldraw[headCol] (\R, \R) circle (0.13);
    \node[right=3pt, font=\small, headCol!60!black] at (\R+0.13, \R+0.25) {\textbf{2}};
    \draw[->, very thick, defscol]
        (\R, \R+0.30) -- (\R-1.1, \R+0.30) node[above, font=\footnotesize] {$N_2$};
    \draw[->, very thick, red!70!black]
        (\R, \R) -- (\R, \R-0.85) node[right, font=\footnotesize] {$m\mathrm{g}$};
    \filldraw[headCol] (0, 2*\R) circle (0.13);
    \node[above=3pt, font=\small, headCol!60!black] at (0, 2*\R+0.13) {\textbf{3}};
    \draw[->, very thick, defscol]
        (-0.22, 2*\R) -- (-0.22, 2*\R-0.85) node[left, font=\footnotesize] {$N_3$};
    \draw[->, very thick, red!70!black]
        (0.22, 2*\R) -- (0.22, 2*\R-0.85) node[right, font=\footnotesize] {$m\mathrm{g}$};
  \end{tikzpicture}%
  }
  \caption{\footnotesize Fart i høyde $h$: $v^2=2\mathrm{g}(H-h)$ (energibevaring fra start).}
  \end{figure}

\end{frame}

%--------------------------------------------------------------------------
%  SLIDE 13 — Løsning: fart i hver posisjon
%--------------------------------------------------------------------------
\begin{frame}{Løsning: fart og normalkraft}

  Energibevaring fra start ($v_0=0$, høyde $H$) til høyde $h$:
  \[
    m\mathrm{g}H = \tfrac{1}{2}mv^2 + m\mathrm{g}h \quad\Longrightarrow\quad v^2 = 2\mathrm{g}(H-h)
  \]

  \pause
  Newtons 2. lov, sentripetalt (mot sentrum):
  \[
    \sum F_{\text{sentripetal}} = \frac{mv^2}{R}
  \]

\end{frame}

%--------------------------------------------------------------------------
%  SLIDE 14 — Løsning: de tre posisjonene
%--------------------------------------------------------------------------
\begin{frame}{}

  \textbf{1 --- bunn} ($N_1$ opp, $m\mathrm{g}$ ned, begge mot/fra sentrum):
  \[
    N_1 - m\mathrm{g} = \frac{mv_1^2}{R} \quad\Longrightarrow\quad N_1 = m\mathrm{g}\!\left(1+\frac{2H}{R}\right) = \doubleunderline{108\,\mathrm{N}}
  \]

  \pause
  \textbf{2 --- siden} ($m\mathrm{g}$ bidrar ikke til sentripetalkraften):
  \[
    N_2 = \frac{mv_2^2}{R} = \frac{2m\mathrm{g}(H-R)}{R} = \doubleunderline{78.5\,\mathrm{N}}
  \]

  \pause
  \textbf{3 --- topp} ($N_3$ og $m\mathrm{g}$ peker begge mot sentrum, nedover):
  \[
    N_3 + m\mathrm{g} = \frac{mv_3^2}{R} \quad\Longrightarrow\quad N_3 = \doubleunderline{49.1\,\mathrm{N}}
  \]

\end{frame}

%--------------------------------------------------------------------------
%  SLIDE 15 — Merk: minste starthøyde
%--------------------------------------------------------------------------
\begin{frame}{}

  Normalkraften er størst i bunnen og minst på toppen.

  \pause
  \merk{Kritisk starthøyde}{
  For at ballen skal holde kontakt med skinna på toppen, må $N_3\geq 0$:
  \[
    \frac{2(H-2R)}{R}\geq 1 \quad\Longrightarrow\quad H \geq \frac{5R}{2}
  \]
  Med $R=2.0\,\mathrm{m}$ er minste tillatte starthøyde $H_{\min}=5.0\,\mathrm{m}$.
  }

\end{frame}

%--------------------------------------------------------------------------
%  SLIDE 16 — Neste gang
%--------------------------------------------------------------------------
\begin{frame}
\begin{center}
  \Huge{Neste gang: effekt og støt}
  \begin{itemize}
    \item Effekt --- arbeid per tidsenhet
    \item Elastisk støt: kollisjon mellom vogner
  \end{itemize}
  \vspace{6mm}
\end{center}
\end{frame}

\end{document}
